\chapter{Building a Vector Database (HNSW \& AST)}

Computing the exact Cosine Similarity for every vector sequentially ($O(N)$) for every user prompt will take seconds or minutes when scaled to large datasets. 

Production Vector DBs (Pinecone, Qdrant) solve this using \textbf{Approximate Nearest Neighbor (ANN)} search algorithms, notably \textbf{HNSW (Hierarchical Navigable Small World)}.

\section{The HNSW Graph Search Algorithm}

HNSW works similarly to a Skip List combined with a social network graph.

\begin{figure}[h]
\centering
\begin{tikzpicture}[
    scale=0.85, transform shape,
    node distance=1.2cm and 0.8cm,
    dot/.style={circle, draw, fill=blue!20, inner sep=2pt, minimum size=0.5cm},
    edge/.style={thick}
]

% Layer 2
\node[dot, fill=red!20] (A) at (0,4) {A};
\node[right=of A, xshift=1.5cm] {Layer 2 (Highest - Entry Point)};

% Layer 1
\node[dot, fill=yellow!20] (B) at (-2,2) {B};
\node[dot, fill=yellow!20] (C) at (2,2) {C};
\node[right=of C, xshift=0.5cm] {Layer 1 (Medium Granularity)};

\draw[edge] (A) -- (B);
\draw[edge] (A) -- (C);

% Layer 0
\node[dot] (D) at (-3,0) {D};
\node[dot] (E) at (-1,0) {E};
\node[dot] (F) at (1,0) {F};
\node[dot] (G) at (3,0) {G};
\node[right=of G, xshift=0.5cm] {Layer 0 (Base Layer)};

\draw[edge] (B) -- (D);
\draw[edge] (B) -- (E);
\draw[edge] (C) -- (F);
\draw[edge] (C) -- (G);
\draw[edge] (D) -- (E);
\draw[edge] (F) -- (G);
\draw[edge] (E) -- (F);

\end{tikzpicture}
\caption{HNSW Graph Search Layers}
\end{figure}

\subsection{Hyperparameters of HNSW}
Configuring HNSW for code search requires tuning specific construction and search hyperparameters:
\begin{itemize}
    \item \textbf{$M$:} The maximum number of bi-directional connections established for each new node during index insertion. Standard values range from 8 to 64. Higher values increase accuracy on high-dimensional vectors (like code) at the cost of VRAM and build time.
    \item \textbf{$M_{\text{max}}$:} The maximum number of connections allowed per node at Layer 0. Typically set to $2M$. If connections exceed this, prune links.
    \item \textbf{$efConstruction$:} The size of the dynamic candidate list evaluated during index construction. A larger list creates a more dense graph, improving search recall but increasing build time.
    \item \textbf{$efSearch$:} The size of the dynamic candidate list evaluated during query routing. Unlike construction parameters, $efSearch$ can be set dynamically at query time to trade off latency for recall.
\end{itemize}

\section{Graph RAG \& Episodic Memory Systems}

Vector similarity searches match isolated content chunks, but fail at holistic relational queries (e.g., ``Find all files that import the `AuthHelper` class and call the `validate` function''). To resolve this, modern runtimes implement **Graph RAG**.

Graph RAG models entities (files, classes, functions) as **Nodes**, and relationships (imports, inheritance, calls) as **Edges**. When querying, the engine traverses relational pathways in a graph database alongside semantic vector matches.


\begin{center}
\noindent\rule{0.6\textwidth}{0.4pt} \\
\vspace{0.3em}
\textit{[Code implementation is available in the companion coding handbook: \texttt{coding-handbook/ch06\_vector\_db/}]} \\
\noindent\rule{0.6\textwidth}{0.4pt}
\end{center}


This graph-based indexing pattern enables the orchestrator to fetch entire dependency hierarchies, preventing compiler syntax crashes when importing modified library code.

\section{Abstract Syntax Tree (AST) Chunking}

Feeding a Vector DB arbitrary string chunks ruins semantic search. If a chunk splits a Python function in half, the vector is garbage.

Production agents parse the \textbf{Abstract Syntax Tree (AST)} using libraries like \texttt{Tree-sitter} to chunk code exactly at the function/class level.

\subsection{AST Chunking Implementation}
Here is a complete, nested-aware Python AST parser that extracts class and function scopes with correct parent references.


\begin{center}
\noindent\rule{0.6\textwidth}{0.4pt} \\
\vspace{0.3em}
\textit{[Code implementation is available in the companion coding handbook: \texttt{coding-handbook/ch06\_vector\_db/}]} \\
\noindent\rule{0.6\textwidth}{0.4pt}
\end{center}

